% ============================================================
%  CAPÍTULO 4 — ESTADO DEL ARTE Y TECNOLOGÍAS
% ============================================================
\chapter{Estado del arte y tecnologías empleadas}
\label{ch:arte}

\section{Desarrollo rápido de aplicaciones web}
\label{sec:rad}

El concepto de \gls{rad} fue formalizado por James Martin en los años noventa~\cite{martin1991rad} como
una respuesta a los ciclos de desarrollo en cascada, demasiado rígidos para adaptarse a los
cambios de requisitos. En el contexto de las aplicaciones web modernas, el \gls{rad} se
manifiesta principalmente a través de tres estrategias:

\begin{enumerate}
  \item \textbf{Scaffolding:} generación automática de código boilerplate a partir de un
    modelo de datos (p.~ej., Django Admin, Rails Scaffolding).
  \item \textbf{Frameworks declarativos:} el desarrollador describe \textit{qué} quiere
    (esquemas, configuraciones) y el \textit{framework} genera \textit{cómo} lograrlo.
  \item \textbf{Low-code / No-code:} plataformas con interfaces visuales que eliminan
    casi por completo la escritura de código (p.~ej., Retool, Appsmith).
\end{enumerate}

El enfoque adoptado en este proyecto se sitúa entre las estrategias 1 y 2: se parte de
modelos de datos explícitos (SQLAlchemy) y se genera automáticamente tanto el
\gls{api} REST como la interfaz de usuario, mediante una tubería de transformación de
esquemas.

\subsection{Comparativa con soluciones existentes}
\label{subsec:comparativa}

\begin{table}[H]
  \centering
  \caption{Comparativa de soluciones RAD web de código abierto.}
  \label{tab:comparativa-rad}
  \begin{tabular}{lcccc}
    \toprule
    \textbf{Solución} & \textbf{Backend} & \textbf{Frontend} & \textbf{Extensible} &
    \textbf{Schema-driven} \\
    \midrule
    Django Admin & Python & Integrado & Limitado & No \\
    Rails Scaffold & Ruby & ERB & Sí & No \\
    Retool & SaaS & Visual & Sí & Parcial \\
    Appsmith & Node.js & React & Sí & Parcial \\
    \textbf{uniback + ngt-gui} & \textbf{Python/FastAPI} & \textbf{Angular} &
    \textbf{Sí} & \textbf{Sí} \\
    \bottomrule
  \end{tabular}
\end{table}

La principal diferencia respecto a las soluciones existentes es que \textit{uniback} +
\textit{ngt-gui} están diseñados para integrarse en proyectos con arquitecturas propias,
sin imponer un stack cerrado ni requerir una plataforma SaaS.

\section{Tecnologías del backend}
\label{sec:tecnologias-backend}

\subsection{Python y FastAPI}
\label{subsec:fastapi}

\gls{fastapi}~\cite{fastapi} es un framework web asíncrono basado en \textbf{ASGI}
(Asynchronous Server Gateway Interface) que utiliza las anotaciones de tipo de Python para
generar automáticamente documentación \textbf{OpenAPI~3.x}. Su principal ventaja respecto
a alternativas como Flask o Django~REST~Framework es la integración nativa con Pydantic
para la validación de datos y la generación de esquemas.

\subsection{SQLAlchemy y Pydantic}
\label{subsec:orm}

\gls{sqlalchemy}~\cite{sqlalchemy} es el \gls{orm} de referencia en el ecosistema Python.
La versión~2.x, utilizada en este proyecto, introduce una \gls{api} basada en
\texttt{mapped\_column} que simplifica la definición de columnas y facilita la introspección
de metadatos, clave para la generación automática de esquemas.

\gls{pydantic}~\cite{pydantic} actúa como capa de validación y serialización entre
SQLAlchemy y el \gls{api} REST. Su método \texttt{model\_json\_schema()} genera un
JSON~Schema estándar a partir de cualquier modelo Pydantic, que este proyecto extiende con
anotaciones \texttt{x-*}.

\subsection{JSON Schema y OpenAPI}
\label{subsec:jsonschema}

JSON~Schema~\cite{jsonschema} es el estándar utilizado como lenguaje común entre el
\textit{backend} y el \textit{frontend}. OpenAPI~3.x~\cite{openapi} lo adopta como
mecanismo de descripción de modelos y permite añadir extensiones de proveedor (campos con
prefijo \texttt{x-}) sin romper la compatibilidad con validadores estándar.

\subsection{PostgreSQL y Docker}
\label{subsec:postgres-docker}

\gls{docker}~\cite{docker} proporciona la infraestructura de contenedores que garantiza la
reproducibilidad del entorno. El servicio de base de datos utilizado es PostgreSQL~16, que
ofrece soporte nativo para columnas \texttt{JSONB}, búsqueda de texto completo
(\texttt{tsvector}) y generación de UUIDs, características aprovechadas por algunos modelos
del proyecto.

\section{Tecnologías del frontend}
\label{sec:tecnologias-frontend}

\subsection{Angular 19}
\label{subsec:angular}

\gls{angular}~\cite{angular} es el framework \gls{spa} elegido por el \gls{itc} para el
desarrollo de sus aplicaciones. La versión~19 introduce mejoras significativas en el
rendimiento (Signals, defer blocks) y consolida el modelo de componentes independientes
(\textit{standalone components}), que se utiliza de forma sistemática en \textit{ngt-gui}.

\subsection{Angular Formly}
\label{subsec:formly}

\gls{formly}~\cite{formly} es la librería que permite generar formularios reactivos de
Angular a partir de una configuración JSON. Su modelo de extensión mediante tipos y
\textit{wrappers} personalizados se adapta perfectamente a la estrategia declarativa del
proyecto: el \textit{backend} emite un JSON~Schema y el \textit{frontend} lo convierte en
una configuración Formly para renderizar el formulario.

\subsection{AG Grid Community}
\label{subsec:aggrid}

\gls{aggrid}~\cite{aggrid} es el componente de tabla de alto rendimiento utilizado para la
cuadrícula de edición masiva. Su edición \textit{inline}, soporte para estados de fila
personalizados y \gls{api} de columnas programáticas lo convierten en la opción idónea para
implementar una interfaz tipo Excel sobre datos dinámicos.

\subsection{NG-ZORRO Antd}
\label{subsec:ngzorro}

NG-ZORRO~\cite{ngzorro} es la librería de componentes \gls{ui} basada en Ant~Design
utilizada para los elementos de interfaz no relacionados con la cuadrícula (menús, botones,
formularios simples, notificaciones). Proporciona un sistema de diseño coherente y
accesible sin imponer dependencias en la lógica de negocio.

\section{Patrones arquitectónicos relacionados}
\label{sec:patrones}

\subsection{Schema-first development}
\label{subsec:schema-first}

El enfoque \textit{schema-first} propone definir el contrato de datos antes de implementar
las capas de negocio y presentación~\cite{schemafirst}. En este proyecto el esquema no se
define manualmente sino que se \textbf{genera automáticamente} desde los modelos ORM,
combinando las ventajas del \textit{schema-first} con las del \textit{model-first}.

\subsection{Plugin pattern}
\label{subsec:plugin-pattern}

El sistema de plugins por capas implementado en \textit{uniback} sigue el patrón
\textbf{Chain of Responsibility}~\cite{gof}: cada plugin recibe el esquema producido por
la capa anterior, lo enriquece y lo pasa a la siguiente. Esto permite añadir o quitar
plugins sin modificar el núcleo del generador.

\subsection{CQRS y operaciones en bloque}
\label{subsec:cqrs}

El endpoint \texttt{/bulk} introduce una separación implícita entre comandos de escritura
(crear, actualizar, eliminar) y consultas de lectura, en la línea del patrón
\gls{crud}~\cite{fowler2002} extendido. El modo \textit{partial} (con savepoints SQL) y
el modo \textit{abort} (transacción global) ofrecen semántica transaccional configurable
por el cliente.
